\section{The fixed proof-plan language}
\label{sec:plan}

DARC compiles one fixed finite plan at a time.  The plan language is broad
enough to organize proofs from discrete mathematics, algebra, bounded number
theory, effective analysis, and finite geometry, but deliberately small enough
that every construct has a direct ordinary-\(\iPRA\) interpretation.

\subsection{Nodes and interfaces}

A plan consists of a finite directed graph \(G=(V,E)\), a public query node
\(q\), and declarations for resources and local proof interfaces.  Nodes have
one of the following roles.

\begin{description}
  \item[Producer.] Establishes an input theorem, modulus, bound, finite net,
  apartness margin, or certificate-existence formula.  A producer may be
  proved, supplied as a hypothesis, or classified as an open residual.

  \item[Compute.] Applies a fixed primitive-recursive transformation and
  exposes its defining equations.  Examples are polynomial evaluation,
  bounded sums, a recurrence step, code projections, and rational arithmetic.

  \item[Verify.] Checks a fixed decidable relation through a characteristic
  bit and an object-level reflection theorem.

  \item[Case or bounded-formula node.] Compiles finite conjunction,
  disjunction, implication under decidability, and bounded quantifier
  alternation to truth bits or strategy certificates.

  \item[Loop.] Encapsulates a bounded iteration with a fixed state type,
  primitive-recursive step, and quantifier-free invariant.

  \item[Relation.] Represents a finite cloud of possible witnesses and a
  compatibility relation between adjacent layers.  It is compiled by backward
  viability.

  \item[Observer.] Maps a possibly nonunique local witness to the public data
  that actually matter.  Its contract includes exact, quotient, or metric
  compatibility.

  \item[Bridge.] Translates a completed ordinary sequent to a client
  presentation and proves comparison lemmas for the client's symbols.
\end{description}

Every node \(v\) has a conclusion formula \(\phi_v\), dependency list
\(\operatorname{dep}(v)\), and local contract
\[
  \bigl(\phi_u: u\in\operatorname{dep}(v)\bigr),\;B_v(d)
  \;\vdash_{\PRAAx}\;\phi_v,
\]
where \(B_v(d)\) is the list of boundary formulas demanded at state \(d\).
A node is active only if its output or one of its resource effects is demanded.
The graph is a proof graph, not an execution graph: a compute node denotes a
proved equation for a fixed primitive-recursive symbol, and a loop node denotes
an induction proof.

\subsection{Resource declarations}

A resource coordinate has a name, kind, finite capacity, and intended theorem
meaning.  Boolean coordinates have capacity one.  A precision coordinate has
capacity chosen for the requested theorem instance.  This finite cap is not a
claim that an analytic object has only finitely many precisions; the theorem
being compiled has a fixed finite target precision, so only the finite interval
needed for that instance appears in the plan.

Typical declarations are:
\begin{center}
\begin{tabularx}{0.95\textwidth}{@{}l l c X@{}}
\toprule
Name & Kind & Capacity & Meaning \\
\midrule
\code{query:value} & query & 1 & approximate value is requested \\
\code{cert:branch} & certificate & 1 & an ordered branch label must escape \\
\code{precision:f} & precision & \(p_{\max}\) & accuracy requested from \(f\) \\
\code{guard:denom} & guard & 1 & positive denominator margin is needed \\
\code{residual:tail} & residual & 1 & tail-modulus producer is not yet supplied \\
\code{backend:foundation} & backend & 1 & Foundation export is requested \\
\bottomrule
\end{tabularx}
\end{center}

A level may encode more than presence.  For example, a guard coordinate can
record one of finitely many margin classes, and a precision coordinate records
an integer level.  The only structural requirement is a finite chain with a
clear monotone meaning: a larger value requests at least as much resource.

\subsection{Transfer rules}

A local transfer summary says what a node needs when one of its outputs is
needed.  The most transparent executable form is a finite family of threshold
rules:
\[
  \bigwedge_{r\in T} d_r\ge a_r
  \quad\Longrightarrow\quad
  d_s\ge b_s\quad(s\in U).
\]
For example:
\begin{align*}
  \Prec(g)\ge p &\Longrightarrow \Prec(f)\ge p+1,\\
  \Cert(\text{ordered root})=1
    &\Longrightarrow \Guard(\text{separation})=1,\\
  \Prec(\text{series})\ge p_0
    &\Longrightarrow \Res(\text{tail modulus})=1.
\end{align*}
The global transfer \(F\) joins all requirements whose thresholds fire.
Threshold rules are monotone by construction: raising the input can only make
more rules fire, and requirements are added by coordinatewise maximum.

More general monotone transfer functions are allowed.  A Lean formalization
should demand a monotonicity proof for each such function.  Threshold rules are
preferred for generated plans because their monotonicity and provenance are
immediate.

\subsection{Plan normalization}

Before demand analysis, a plan is normalized as follows.
\begin{enumerate}
  \item Inline purely definitional aliases so each compute node has one fixed
        symbol and explicit equations.
  \item Replace each fixed bounded formula by a node that will generate its
        truth bit and reflection derivation.
  \item Isolate runtime cycles as explicit loop nodes with state and invariant.
        A cyclic dependency in precision or liveness is not a runtime loop; it
        remains in the monotone transfer network.
  \item Group mutually recursive finite witness choices into one relation node
        rather than choosing each witness independently.
  \item Attach every partial operation to a guard formula.
  \item Attach every unproved quantitative input to a producer node and a
        residual coordinate.
  \item Compute the strongly connected components of the remaining proof
        graph.  After loop extraction, the condensation graph used for
        simultaneous-cut assembly must be acyclic.
\end{enumerate}

The normalization does not prove local mathematics.  It makes proof
obligations explicit and puts them in forms supported by the checked library.

\subsection{Machine-readable representation}

The accompanying schema \pathcode{schema/darc_plan.schema.json} records:
resource names and capacities, threshold rules, node kinds, dependency lists,
boundary formulas, certificate scope, Lean target, and status.  The schema is
not part of the trusted proof kernel.  Its purpose is reproducibility: a
reviewer can compare the human theorem, demand analysis, generated proof
schedule, and final residual list against one common plan.

A trustworthy generator should emit at least four artifacts:
\begin{enumerate}
  \item the closed demand state and its iteration history;
  \item the pruned node schedule;
  \item a theorem ledger mapping each node to a Lean declaration or residual;
  \item the final ordinary sequent, before any backend translation.
\end{enumerate}
The reference Python code emits the first item and small demonstrations of the
finite strategy, trace, and viability algorithms.  It intentionally does not
parse or accept Lean proof certificates.
